\section{Christ, Mary, and the angels}

\subsection{One Head for the heavenly and earthly Church}

The angels who minister to Jesus serve the incarnate Son. His assumption of
human nature does not place him beneath them in the order of grace. Thomas
asks the stronger question: is Christ \emph{as man} their Head? His answer
begins with a common end. Angels and human beings are called to the enjoyment
of God; their different natures and offices belong to one community ordered
toward that end. Christ possesses the fullness from which this community
receives its life. The unity of its Head is therefore greater than the
diversity of its members.\footnote{\ST{III, q.~8, a.~4, corpus and ad~1--3};
English Dominican translation, revised 1920, New Advent witness,
\url{https://www.newadvent.org/summa/4008.htm}.}

The objection from Hebrews 2:16 is substantial: the Word assumed human nature,
not an angelic nature. Thomas does not answer by giving Christ a third nature.
He answers that human souls and angels agree in intellectual life, although
they differ specifically. Nor does the fact that angels already behold God
remove them from Christ's Church. The earthly Church lives by faith; the
heavenly Church possesses the vision toward which faith journeys. Christ's
headship embraces both. His humanity acts in this spiritual order through its
personal union with the Word. This is a claim about the one Christ and his
gracious influence, not about a human creature exercising independent divine
power.

\begin{quote}\small
\textbf{Summa:} III, q.~8, a.~4.\quad
\textbf{Grade:} Church teaching concerning Christ's supremacy; Thomistic
explanation of his headship according to his humanity.\quad
\textbf{Authorities:} Colossians 2:10; Ephesians 1:20--23; Matthew 4:11.
\end{quote}

This headship gives angelic ministry its Christological direction. An angel
does not provide a route to God that passes around the Son. The minister who
announces his birth, strengthens him in his human suffering, or proclaims his
resurrection serves the same Lord whose glory exceeds the minister's own.
Service does not imply that the Son lacks divine power. It displays the
economy in which his true human life receives created assistance and his
creatures participate in his work. The direction of Christian attention
follows the direction of their service: toward Christ.

\subsection{Gabriel and the Virgin's consent}

The Annunciation joins heavenly ministry to human freedom. In Thomas's account,
Mary hears and believes before she conceives; she becomes a conscious witness
to the mystery accomplished in her. Her consent expresses the acceptance of
human nature in its spiritual marriage with the Word. Predestination does not
make that consent superfluous. God's purpose includes the willing response of
the person whom he prepares for it.\footnote{\ST{III, q.~30, a.~1, corpus and
ad~1--2}; \url{https://www.newadvent.org/summa/4030.htm}.}

Gabriel's embassy is fitting because God communicates divine things through
angels, and because the restoration of humanity answers the earlier deception
of the woman by the serpent. Yet the order of ministry is not an absolute
ranking of the persons involved. Thomas says that Mary is above the angels
in the dignity for which God chose her, while receiving angelic instruction
as a wayfarer. A person may receive a service from someone whose ultimate
dignity is lower than her own. The Annunciation therefore does not diminish
Mary's preeminence. It reveals the cooperation of gifts.\footnote{\ST{III,
q.~30, a.~2, corpus and ad~1}. The argument is Thomas's; the patristic
quotations embedded in this article are not independently attributed here.}

Thomas also resists the inference that the greatest embassy must have been
carried by a seraph. Gabriel is an archangel; Thomas considers it credible
that he is the highest of the archangels, without making him a member of the
highest order. This is a theological judgment about the fittingness of his
office.\footnote{\ST{III,
q.~30, a.~2, ad~4}.}

The bodily appearance of the messenger belongs to the mystery's embodied
character. An invisible minister becomes visible to announce that the
invisible God will take flesh. Mary receives illumination of mind together
with an encounter accessible to her senses. Thomas consequently does not
identify seeing Gabriel's assumed appearance with seeing the angelic essence.
The message proceeds from salutation to instruction and then to consent;
Mary's question seeks the manner of the promised conception, rather than
refusing belief.\footnote{\ST{III, q.~30, aa.~3--4, especially a.~3,
corpus and ad~1, and a.~4, corpus and ad~2}.}

\begin{quote}\small
\textbf{Summa:} III, q.~30, aa.~1--4.\quad
\textbf{Grade:} Church teaching concerning the Annunciation and Mary's free
consent; Thomistic arguments for the fittingness of its manner.\quad
\textbf{Authorities:} Luke 1:26--38; the article's received Dionysian and
patristic authorities.\quad
\textbf{Contrary opinion:} Gabriel's membership in the highest order is
reported and rejected in a.~2, ad~4.
\end{quote}

\subsection{Mary above the angels}

Vatican II teaches that the grace of divine motherhood places Mary above all
other creatures and explicitly describes her exaltation above angels and
human beings. It also identifies her as redeemed through her Son, a member
of the Church, and the Lord's handmaid. Her queenship is received from him;
her veneration differs essentially from the adoration offered to God.
These affirmations belong together. Mary's eminence is the triumph of grace
in a creature, never a second divine sovereignty.\footnote{Second Vatican
Council, \emph{Lumen gentium} (21 November 1964), nn.~53, 56, 59, 66;
official English text, \url{https://www.vatican.va/archive/hist_councils/ii_vatican_council/documents/vat-ii_const_19641121_lumen-gentium_en.html}.}

The distinction between nature and grace becomes especially clear here.
Being human does not become being angelic through sanctification. Mary
remains the mother who gave the Word his human flesh. The height of a created
nature does not measure all that God may give through union with Christ.
Angelic excellence and Marian preeminence can therefore be affirmed without
confusion: one concerns the kind of creature, the other the singular gift and
vocation bestowed upon her.
